\section{Laborexperimente und Präzisionstests}

Während astrophysikalische Beobachtungen die strengsten Grenzen für eine 
mögliche Photonenmasse liefern, spielen Laborexperimente eine wichtige 
komplementäre Rolle. Sie ermöglichen kontrollierte, modellunabhängige Tests 
der Elektrodynamik unter gut definierten Bedingungen. Auch wenn die 
experimentelle Sensitivität im Labor viele Größenordnungen unter den 
astrophysikalischen Grenzen liegt, besitzen diese Tests einen hohen 
konzeptionellen Wert: Sie überprüfen die Grundstruktur der 
Maxwell-Theorie, ohne dass Unsicherheiten über kosmische Magnetfelder, 
Plasmaeffekte oder großskalige Strukturen eine Rolle spielen.

\subsection*{Tests des Coulomb-Gesetzes}

Eine der direktesten Folgen einer von Null verschiedenen Photonenmasse ist 
die Modifikation des Coulomb-Potentials. Die Proca-Theorie sagt ein 
Yukawa-artiges Potential voraus,
\[
V(r) = \frac{q}{4\pi\varepsilon_0} \frac{e^{-m_\gamma r}}{r},
\]
anstelle des klassischen $1/r$-Verhaltens. Präzise Messungen der 
elektrostatischen Kraft setzen daher Grenzen an $m_\gamma$.

Versuche mit Torsionswaagen, elektrostatischen Nullmessungen und 
hochpräzisen Kondensatoranordnungen ergeben typischerweise
\[
m_\gamma \lesssim 10^{-14}\,\text{eV}.
\]
Dieser Wert ist deutlich schwächer als astrophysikalische Grenzen, beruht 
jedoch ausschließlich auf Laborphysik und minimalen theoretischen 
Voraussetzungen.

\subsection*{Messungen magnetischer Felder}

Ein massives Photon würde bewirken, dass statische Magnetfelder nur eine 
endliche Reichweite besitzen. Experimente mit langen Solenoiden, 
supraleitender magnetischer Abschirmung und Präzisionsmagnetometrie wurden 
genutzt, um solche Effekte zu testen.

Messungen des Magnetfeldes in supraleitenden Kavitäten und um ausgedehnte 
Solenoidstrukturen führen zu Einschränkungen der Größenordnung
\[
m_\gamma \lesssim 10^{-13}\,\text{eV}.
\]
Diese Tests gelten als besonders zuverlässig, da die Magnetfeldgeometrie 
mit hoher Genauigkeit berechenbar ist.

\subsection*{Resonator-Experimente}

Resonante elektromagnetische Kavitäten bieten eine weitere hochempfindliche 
Methode, Abweichungen von den Maxwell-Gleichungen zu untersuchen. Eine 
Photonenmasse würde Resonanzfrequenzen verschieben und die Modenstruktur 
verändern.

Der Vergleich zwischen gemessenen und theoretisch vorhergesagten 
Resonanzspektren ergibt typischerweise
\[
m_\gamma \lesssim 10^{-12}\,\text{eV}.
\]
Diese Experimente sind technologisch anspruchsvoll, testen aber die 
Wellennatur des elektromagnetischen Feldes in einer äußerst präzisen, 
kontrollierten Umgebung.

\subsection*{Interferometrie und Ausbreitungstests}

Interferometrische Verfahren können winzige Abweichungen in der 
Wellenfortpflanzung detektieren. Trotz kurzer Laborbaselines erreichen 
moderne Interferometer eine Sensitivität, die strenge Nulltests des 
Maxwell-Verhaltens ermöglicht.

Experimente, die einsetzen:
\begin{itemize}
	\item Mikrowelleninterferometrie,
	\item optische Verzögerungslinien,
	\item Laser-Resonatoren,
\end{itemize}
haben nach masseninduzierten Dispersions- oder Phasenanomalien gesucht.  
Alle Ergebnisse sind konsistent mit der Maxwell-Theorie und führen zu 
Grenzen der Größenordnung
\[
m_\gamma \lesssim 10^{-13}\,\text{eV}.
\]

\subsection*{Vergleich mit astrophysikalischen Grenzen}

Labortests können in ihrer absoluten Sensitivität nicht mit 
astrophysikalischen Grenzen konkurrieren. Dennoch bilden sie eine 
unverzichtbare Grundlage:

\begin{itemize}
	\item Sie beruhen auf einfachen, gut kontrollierbaren Geometrien.
	\item Sie sind frei von Annahmen über kosmische Magnetfelder oder 
	Plasmadichten.
	\item Sie testen elektromagnetische Strukturen unter reproduzierbaren 
	Bedingungen.
	\item Ihre Nullresultate stimmen vollständig mit den Vorhersagen eines 
	masselosen Photons überein.
\end{itemize}

Die Übereinstimmung zwischen Labor- und Astrophysik – trotz völlig 
verschiedener physikalischer Regime – stärkt die Schlussfolgerung, dass 
das Photon mit extrem hoher Genauigkeit masselos ist.

\subsection*{Übergang zu konzeptionellen Argumenten aus der QED}

Das nächste Kapitel führt von den empirischen Tests zu den 
theoretisch-konzeptionellen Überlegungen. In der 
Quantenelektrodynamik ist die Masselosigkeit des Photons eng mit 
Eichsymmetrie, Renormierbarkeit und der Struktur quantenmechanischer 
Amplituden verknüpft. Diese Aspekte bilden die tiefsten Gründe dafür, 
warum eine Photonenmasse mit dem Standardmodell unvereinbar ist.

Laborexperimente vervollständigen damit die empirische Grundlage des 
Gesamtarguments: Falls das Photon eine Masse besitzt, muss sie weit unter 
halb der Nachweisgrenze heutiger Experimente liegen.
